% 릴레이식 기계의 개략 치수(앞, 옆). 단위 mm. 그림 fig:hw-relay-dim-machine
% 근거: hardware/relay/layout.md 3.5, 3.7, 3.11, 3.12절. 프레임 1,880 = 40 + 6 x 240 + 360 + 40,
% 뱅크 835 = 60 + 175(로터 띠) + 290 x 2 + 20 (카트리지식, 2026-11-16; 카드식은 660).
% 로터 띠 930-1,105: 유닛마다 폭 120에 카트리지 자리 셋(피치 39), 카트리지 약 958-1,077 (common/rotor-cartridge.md 3.3절).
\begin{tikzpicture}[x=0.0052cm,y=0.0052cm,
    dim/.style={<->,cgray,line width=0.4pt},
    dl/.style={font=\tiny,text=black,inner sep=1pt},
    lab/.style={font=\tiny,align=center}]
  % ================= (가) 앞 =================
  \node[font=\footnotesize\bfseries,anchor=west] at (0,1860) {(가) 앞};
  \draw[cgray,line width=0.6pt] (-330,0) -- (2620,0);
  % 받침대
  \draw[draw=cgray,dashed] (20,0) rectangle (1860,400);
  \node[lab,text=cgray] at (940,200) {받침대};
  % 패널
  \draw[draw=cgray,fill=white,line width=0.6pt] (0,400) rectangle (1880,870);
  \node[lab] at (760,610) {공통부 패널 (기어식과 같은 설계): 글자 잭, 대각선 보드, 테스트 레지스터};
  \foreach \i in {0,...,11} {
    \draw[draw=cgray,fill=clight] ({40+\i*120+12},790) rectangle ++(42,36);
    \draw[draw=cgray,fill=clight] ({40+\i*120+66},790) rectangle ++(42,36);
  }
  \node[dl,anchor=west] at (40,755) {유닛 잭 12쌍, 피치 120};
  \draw[draw=cgray] (1480,450) rectangle (1840,870);
  \node[lab] at (1660,640) {조작부\\비상 정지};
  % 뱅크 틀
  \draw[draw=cgray,fill=clight,line width=0.6pt] (0,870) rectangle (1880,1705);
  % 다이얼 띠
  \draw[draw=cgray,fill=white] (40,870) rectangle (1840,930);
  \foreach \i in {0,...,11} { \draw[draw=cgray,fill=cyellowbg] ({40+\i*120+60},900) circle (20); }
  % 로터 띠: 유닛마다 카트리지 자리 셋 (L, M, R, 피치 39). 1-6번은 빈 자리, 7-12번은 예로 II, IV, V
  \draw[draw=cgray,fill=white] (40,930) rectangle (1480,1105);
  \foreach \i in {0,...,11} {
    \draw[draw=cgray,line width=0.3pt] ({40+\i*120},930) -- ({40+\i*120},1105);
  }
  \foreach \i in {0,...,5} {
    \foreach \b in {-1,0,1} {
      \draw[draw=cgray,dashed,line width=0.3pt] ({40+\i*120+60+\b*39-15},958) rectangle ++(30,119);
    }
  }
  \foreach \i in {6,...,11} {
    \foreach \b/\col in {-1/brown!75!red, 0/yellow!85!orange, 1/brown!35!white} {
      \draw[draw=black!60,fill=\col,line width=0.3pt] ({40+\i*120+60+\b*39-15},958) rectangle ++(30,119);
    }
  }
  % 유닛 자리: 1-3열은 빈 자리, 4-6열은 채움
  \foreach \c in {0,1,2} {
    \foreach \r in {0,1} {
      \draw[draw=cgray,dashed,fill=white] ({40+\c*240},{1105+\r*290}) rectangle ++(240,290);
      \pgfmathtruncatemacro{\n}{2*\c+1+\r}
      \node[lab,text=cgray] at ({40+\c*240+120},{1105+\r*290+145}) {\n};
    }
  }
  \foreach \c in {3,4,5} {
    \foreach \r in {0,1} {
      \draw[draw=cgray,fill=white] ({40+\c*240},{1105+\r*290}) rectangle ++(240,290);
      \draw[draw=cblue,fill=cbluebg] ({40+\c*240+6},{1105+\r*290+20}) rectangle ++(228,250);
      \foreach \b in {0,...,4} {
        \draw[cblue,line width=1.2pt] ({40+\c*240+18},{1105+\r*290+58+\b*38}) -- ++(204,0);
      }
      \pgfmathtruncatemacro{\n}{2*\c+1+\r}
      \node[lab,fill=cbluebg,inner sep=1pt] at ({40+\c*240+120},{1105+\r*290+153}) {\n};
    }
  }
  % 제어 칸 (보드 자리는 유닛과 함께 175 올림)
  \draw[draw=cgray,fill=white] (1480,930) rectangle (1840,1685);
  \draw[draw=cgreen,fill=cgreenbg] (1487,1425) rectangle (1737,1653);
  \node[lab] at (1612,1539) {공통\\계수기};
  \draw[draw=cgray,fill=cyellowbg] (1487,1135) rectangle (1737,1363);
  \node[lab] at (1612,1249) {순서기};
  % 위치 표시 보드 88 x 294 (글자 세 줄 l, m, r)
  \draw[draw=corange,fill=corangebg] (1746,1255) rectangle (1834,1549);
  \foreach \j in {0,1,2} {
    \foreach \k in {0,...,25} {
      \fill[cgray!60] ({1760+\j*29},{1270+\k*10.5}) circle (2.5);
    }
  }
  \node[lab] at (1790,1205) {위치\\표시};
  \node[lab,text=corange] at (1660,1017) {$\leftarrow$ 로터 띠\\(카트리지 3 × 12)};
  % 화면
  \draw[draw=cgray,dashed] (1540,1710) rectangle (1780,1865);
  \node[dl] at (1660,1787) {화면};
  % 손이 닿는 범위
  \draw[cred,line width=0.4pt,dashed] (-330,380) -- (0,380);
  \draw[cred,line width=0.4pt,dashed] (-330,1220) -- (1880,1220);
  \node[dl,text=cred,anchor=south west] at (-330,385) {380};
  \node[dl,text=cred,anchor=south west] at (-330,1225) {1,220};
  % 치수: 가로
  \foreach \x in {0,40,1480,1840,1880} { \draw[cgray,line width=0.3pt] (\x,-10) -- (\x,-150); }
  \draw[dim] (40,-110) -- (1480,-110) node[midway,dl,fill=white] {6 × 240 = 1,440};
  \draw[dim] (1480,-110) -- (1840,-110) node[midway,dl,fill=white] {360};
  \node[dl] at (20,-60) {40};
  \node[dl] at (1860,-60) {40};
  \draw[cgray,line width=0.3pt] (0,-150) -- (0,-290);
  \draw[cgray,line width=0.3pt] (1880,-150) -- (1880,-290);
  \draw[dim] (0,-250) -- (1880,-250) node[midway,dl,fill=white] {1,880};
  % 치수: 세로 (왼쪽)
  \foreach \h in {0,400,870,930,1105,1395,1685,1705} { \draw[cgray,line width=0.3pt] (-10,\h) -- (-150,\h); }
  \draw[dim] (-110,0) -- (-110,400) node[midway,dl,fill=white] {400};
  \draw[dim] (-110,400) -- (-110,870) node[midway,dl,fill=white] {470};
  \draw[dim] (-110,930) -- (-110,1105) node[midway,dl,fill=white] {175};
  \draw[dim] (-110,1105) -- (-110,1395) node[midway,dl,fill=white] {290};
  \draw[dim] (-110,1395) -- (-110,1685) node[midway,dl,fill=white] {290};
  \node[dl,anchor=east] at (-20,900) {60};
  \node[dl,anchor=east] at (-20,1695) {20};
  \draw[cgray,line width=0.3pt] (-150,0) -- (-270,0);
  \draw[cgray,line width=0.3pt] (-150,1705) -- (-270,1705);
  \draw[dim] (-230,0) -- (-230,1705) node[midway,dl,fill=white,rotate=90] {1,705};
  % ================= (나) 옆 =================
  \begin{scope}[shift={(2080,0)}]
    \node[font=\footnotesize\bfseries,anchor=west] at (-40,1860) {(나) 옆 (공기, 리드)};
    \draw[draw=cgray,dashed] (20,0) rectangle (380,400);
    \draw[draw=cgray,fill=white,line width=0.6pt] (0,400) rectangle (400,870);
    \node[lab] at (170,560) {패널};
    \draw[draw=cgray,dashed] (240,440) rectangle (390,840);
    \node[lab,text=cgray] at (315,690) {전원부};
    \draw[draw=cgray,fill=clight,line width=0.6pt] (70,870) rectangle (400,1705);
    % 유닛 앞 덮개 (로터 띠 위만)
    \draw[cblue,line width=0.8pt] (73,1105) -- (73,1685);
    % 카드 겹침: 로터 카드의 자리(150, 200, 250)는 빈다
    \foreach \r in {0,1} {
      \foreach \z in {100,125,175,225,275,315} {
        \draw[cblue,line width=0.5pt] (\z,{1105+\r*290+20}) -- (\z,{1105+\r*290+270});
      }
      \foreach \z in {150,200,250} {
        \draw[cgray,dashed,line width=0.3pt] (\z,{1105+\r*290+20}) -- (\z,{1105+\r*290+270});
      }
    }
    \draw[cgray,line width=0.5pt] (328,930) -- (328,1665);
    % 로터 띠: 틀, 받침 기판, 뒷면 D-sub, 앞의 카트리지
    \draw[draw=cgray,fill=clight] (78,940) rectangle (84,1095);
    \draw[cgreen,line width=0.8pt] (86,938) -- (86,1097);
    \draw[draw=cgray,fill=white] (88,960) rectangle (100,1000);
    \draw[draw=cgray,fill=white] (88,1035) rectangle (100,1075);
    \draw[draw=black!60,fill=yellow!85!orange] (63,958) rectangle (77,1077);
    \node[dl,anchor=east] at (58,1017) {카트리지};
    % 리드: 겹침의 빈 자리 옆 시프터에서 뒤 배선 공간으로 내려와 받침 기판 뒤로
    \draw[corange,line width=0.8pt] (225,1105) -- (225,1090) -- (360,1090);
    \draw[corange,line width=0.8pt] (225,1395) -- (225,1385) -- (360,1385);
    \draw[corange,line width=0.8pt] (360,1385) -- (360,1040) -- (100,1040);
    \node[dl,anchor=west,text=corange] at (405,1200) {리드};
    % 팬과 흡기구
    \draw[draw=cgray,fill=cyellowbg] (372,1570) rectangle (397,1680);
    \node[dl,anchor=west] at (405,1585) {팬};
    \node[dl,anchor=west] at (405,930) {흡기 (필터)};
    % 공기
    \draw[->,cred,line width=0.8pt] (470,880) -- (340,880) -- (130,915);
    \foreach \z in {137,212,262} { \draw[->,cred,line width=0.6pt] (\z,935) -- (\z,1675); }
    \draw[->,cred,line width=0.8pt] (200,1693) -- (350,1693) -- (372,1645) -- (470,1645);
    % 깊이
    \draw[cgray,line width=0.3pt] (0,-10) -- (0,-150);
    \draw[cgray,line width=0.3pt] (400,-10) -- (400,-150);
    \draw[dim] (0,-110) -- (400,-110) node[midway,dl,fill=white] {400};
    \draw[cgray,line width=0.3pt] (70,1710) -- (70,1805);
    \draw[cgray,line width=0.3pt] (400,1710) -- (400,1805);
    \draw[dim] (70,1765) -- (400,1765) node[midway,dl,fill=white] {330};
  \end{scope}
\end{tikzpicture}
